\title{Επισκόπηση Λειτουργίας Συστήματος Μεταφοράς Ηλεκτρικής Ενέργειας}

\section{Θεωρία}
\label{kef-01-theoria}

\subsection{Ρόλος και απαιτήσεις λειτουργίας του Συστήματος Μεταφοράς}

Το Σύστημα Μεταφοράς Ηλεκτρικής Ενέργειας (ΣΜΗΕ) είναι το δίκτυο υπερυψηλής και υψηλής τάσης
που διασυνδέει τους σταθμούς παραγωγής με τα κέντρα κατανάλωσης και με τα γειτονικά συστήματα.
Στην Ελλάδα λειτουργεί στα επίπεδα των 400 kV, 150 kV και 66 kV και τη διαχείρισή του ασκεί
ο Ανεξάρτητος Διαχειριστής Μεταφοράς Ηλεκτρικής Ενέργειας (ΑΔΜΗΕ), ο οποίος είναι μέλος του
ENTSO-E και τμήμα της σύγχρονης περιοχής της Ηπειρωτικής Ευρώπης (Continental Europe) [15].

Ο Διαχειριστής Συστήματος Μεταφοράς (Transmission System Operator, TSO) καλείται να ικανοποιεί
ταυτόχρονα τρεις απαιτήσεις, οι οποίες διατρέχουν ολόκληρο το παρόν κεφάλαιο:

\begin{itemize}
\item \textbf{Ισοζύγιο ισχύος.} Η ηλεκτρική ενέργεια δεν αποθηκεύεται σε σημαντικές ποσότητες
      εντός του συστήματος. Παραγωγή και ζήτηση πρέπει να εξισώνονται συνεχώς, διαφορετικά η
      διαφορά καλύπτεται από την κινητική ενέργεια των στρεφόμενων μαζών και εκδηλώνεται ως
      μεταβολή της συχνότητας [8].
\item \textbf{Ασφάλεια λειτουργίας.} Το σύστημα πρέπει να παραμένει εντός των ορίων ασφαλούς
      λειτουργίας (θερμικά όρια στοιχείων, όρια τάσεων, όρια ευστάθειας) όχι μόνο στην τρέχουσα
      κατάσταση, αλλά και μετά από κάθε πιθανή μεμονωμένη διαταραχή.
\item \textbf{Ποιότητα.} Η συχνότητα και οι τάσεις πρέπει να διατηρούνται εντός προδιαγεγραμμένων
      περιθωρίων, τα οποία για την Ηπειρωτική Ευρώπη ορίζονται ρυθμιστικά [12, 14].
\end{itemize}

\subsubsection*{Καταστάσεις λειτουργίας}

Ο Κανονισμός (ΕΕ) 2017/1485 (System Operation Guideline, SO GL) ταξινομεί στο Άρθρο 18 τη
λειτουργία του συστήματος σε πέντε καταστάσεις [12]:

\begin{enumerate}
\item \textbf{Κανονική} (normal): όλα τα μεγέθη εντός ορίων και το κριτήριο N-1 ικανοποιείται.
\item \textbf{Επιφυλακής} (alert): τα μεγέθη παραμένουν εντός ορίων, αλλά έχει εξαντληθεί το
      περιθώριο εφεδρείας ή παραβιάζεται το N-1.
\item \textbf{Έκτακτης ανάγκης} (emergency): παραβιάζεται τουλάχιστον ένα όριο ασφαλούς λειτουργίας.
\item \textbf{Συσκότισης} (blackout): εκτεταμένη απώλεια φορτίου ή κατάρρευση τμήματος του συστήματος.
\item \textbf{Αποκατάστασης} (restoration): ο Διαχειριστής έχει ενεργοποιήσει το σχέδιο επανατροφοδότησης.
\end{enumerate}

Η ταξινόμηση αυτή ανάγεται στην κλασική εργασία του Dy Liacco [19], η οποία εισήγαγε την ιδέα ότι
ο έλεγχος του συστήματος οφείλει να προσαρμόζεται στην εκάστοτε κατάσταση ασφαλείας — από τον
προληπτικό έλεγχο στην κανονική κατάσταση έως τον διορθωτικό και τον έλεγχο αποκατάστασης.
Το \textbf{κριτήριο N-1} αποτελεί τον πρακτικό ορισμό της ασφάλειας: το σύστημα πρέπει να μπορεί να
απορροφήσει την απώλεια οποιουδήποτε μεμονωμένου στοιχείου (γραμμής, μετασχηματιστή, μονάδας
παραγωγής) χωρίς παραβίαση ορίων και χωρίς αλυσιδωτές αποσυνδέσεις.

Ο έλεγχος αυτός προϋποθέτει, όμως, ότι ο Διαχειριστής \emph{γνωρίζει} την κατάσταση του συστήματος.
Οι επόμενες ενότητες εξετάζουν πρώτα την υποδομή που παράγει τις μετρήσεις, στη συνέχεια τον
αλγόριθμο που τις μετατρέπει σε αξιόπιστη εικόνα του δικτύου, και τέλος τις χρονικές κλίμακες στις
οποίες λαμβάνονται οι αποφάσεις.

\subsection{Εποπτεία και έλεγχος: SCADA και EMS}

Η βασική υποδομή τηλεμέτρησης και τηλεχειρισμού του ΣΜΗΕ είναι το σύστημα \textbf{SCADA}
(Supervisory Control and Data Acquisition). Στους υποσταθμούς εγκαθίστανται Απομακρυσμένες
Τερματικές Μονάδες (Remote Terminal Units, RTU) και ευφυείς ηλεκτρονικές συσκευές (Intelligent
Electronic Devices, IED), οι οποίες συλλέγουν:

\begin{itemize}
\item \textbf{Αναλογικές μετρήσεις}: ενεργός και άεργος ροή ισχύος στα άκρα των γραμμών και των
      μετασχηματιστών, εγχεόμενη ισχύς σε ζυγούς, μέτρα τάσεων, ρεύματα κλάδων.
\item \textbf{Ψηφιακές ενδείξεις}: θέσεις διακοπτών και αποζευκτών, θέσεις λήψεων μετασχηματιστών,
      σήματα συναγερμών και λειτουργίας προστασιών.
\end{itemize}

Τα δεδομένα μεταφέρονται στο Κέντρο Ελέγχου Ενέργειας μέσω τυποποιημένων πρωτοκόλλων τηλεπροστασίας
και τηλεμετρίας, με πλέον διαδεδομένα τα IEC 60870-5-101/104, το DNP3 και — για την επικοινωνία
εντός του υποσταθμού — το πρότυπο IEC 61850 [18]. Ο ρυθμός σάρωσης ενός συμβατικού SCADA είναι
τυπικά της τάξης των λίγων δευτερολέπτων (συνήθως 2–10 s).

Δύο χαρακτηριστικά του SCADA είναι κρίσιμα για όσα ακολουθούν:

\begin{itemize}
\item Οι μετρήσεις \textbf{δεν είναι χρονικά συγχρονισμένες} μεταξύ τους. Συλλέγονται σε
      διαφορετικές στιγμές εντός του κύκλου σάρωσης, οπότε το σύνολο των μετρήσεων αποτελεί μια
      \emph{ημι-στατική} απεικόνιση (snapshot) και όχι στιγμιότυπο.
\item Οι μετρήσεις είναι \textbf{θορυβώδεις και ενίοτε εσφαλμένες}, ενώ ορισμένα μεγέθη
      (π.χ.\ γωνίες τάσης) δεν μετρώνται καθόλου από το συμβατικό SCADA.
\end{itemize}

Το \textbf{EMS} (Energy Management System) είναι το λογισμικό που αξιοποιεί αυτά τα δεδομένα.
Οι κύριες λειτουργίες του, με τη σειρά που εκτελούνται, είναι ο επεξεργαστής τοπολογίας
(topology processor), η \emph{εκτίμηση κατάστασης}, η ανάλυση ενδεχομένων (contingency analysis),
η βέλτιστη ροή ισχύος (Optimal Power Flow) και ο αυτόματος έλεγχος παραγωγής (Automatic Generation
Control, AGC) [6, 7].

\subsection{Συγχρονισμένες μετρήσεις φασιθετών (PMU) και συστήματα WAMS}

Οι \textbf{Μονάδες Μέτρησης Φασιθετών} (Phasor Measurement Units, PMU) αίρουν τον βασικό
περιορισμό του SCADA. Μια PMU υπολογίζει τον φασιθέτη (μέτρο και \emph{γωνία}) της τάσης και του
ρεύματος και τον συνοδεύει με χρονοσήμανση από κοινή χρονική αναφορά UTC, τυπικά μέσω GPS. Επειδή
όλες οι PMU αναφέρονται στην ίδια χρονική βάση, οι μετρήσεις τους είναι άμεσα συγκρίσιμες σε
ολόκληρο το σύστημα — εξ ου και ο όρος \emph{συγχρονοφασιθέτες} (synchrophasors) [9].

Τα βασικά τεχνικά χαρακτηριστικά τους καθορίζονται από το πρότυπο IEC/IEEE 60255-118-1:2018 [11],
το οποίο διαδέχθηκε τη σειρά IEEE C37.118:

\begin{itemize}
\item \textbf{Ρυθμός αναφοράς}: τυπικά 10, 25 ή 50 πλαίσια ανά δευτερόλεπτο για σύστημα 50 Hz,
      με δυνατότητα υψηλότερων ρυθμών — δηλαδή δύο έως τρεις τάξεις μεγέθους ταχύτερα από το SCADA.
\item \textbf{Ακρίβεια}: το συνολικό σφάλμα φασιθέτη (Total Vector Error, TVE) δεν πρέπει να
      υπερβαίνει το 1\% στη μόνιμη κατάσταση. Ορίζονται επιπλέον το σφάλμα συχνότητας (FE) και
      το σφάλμα ρυθμού μεταβολής συχνότητας (RFE).
\item \textbf{Κλάσεις}: η κλάση P (protection) βελτιστοποιείται για ταχεία απόκριση και η κλάση M
      (measurement) για ακρίβεια υπό παρουσία αρμονικών και παρεμβολών.
\end{itemize}

Οι μετρήσεις συγκεντρώνονται ιεραρχικά σε Συμπυκνωτές Δεδομένων Φασιθετών (Phasor Data
Concentrators, PDC) και συνθέτουν ένα \textbf{Σύστημα Εποπτείας Ευρείας Περιοχής} (Wide Area
Monitoring System, WAMS). Οι χαρακτηριστικές εφαρμογές είναι η παρακολούθηση ηλεκτρομηχανικών
ταλαντώσεων μεταξύ περιοχών, η μέτρηση του ρυθμού μεταβολής της συχνότητας, η ανίχνευση
νησιδοποίησης, η επικύρωση μοντέλων μετά από διαταραχές και — όπως θα δούμε παρακάτω — η ενίσχυση
ή και η αντικατάσταση της κλασικής εκτίμησης κατάστασης [9].

\subsection{Διατάξεις ελεγχόμενης μεταφοράς: FACTS}

Ενώ SCADA και PMU \emph{παρατηρούν} το σύστημα, οι διατάξεις \textbf{FACTS} (Flexible AC Transmission
Systems) το \emph{ελέγχουν} ενεργά. Πρόκειται για διατάξεις ηλεκτρονικών ισχύος που επιτρέπουν την
ταχεία και συνεχή ρύθμιση των παραμέτρων που καθορίζουν τη ροή ισχύος σε ένα δίκτυο εναλλασσομένου
ρεύματος — της τάσης, της σύνθετης αντίστασης και της γωνίας [10]. Κατατάσσονται ανάλογα με τον
τρόπο σύνδεσής τους:

\begin{itemize}
\item \textbf{Παράλληλης σύνδεσης}: ο Στατός Αντισταθμιστής Αέργου Ισχύος (Static Var Compensator,
      SVC) και ο STATCOM. Ελέγχουν την τάση του ζυγού μέσω έγχυσης ή απορρόφησης αέργου ισχύος και
      συνεισφέρουν στη δυναμική στήριξη τάσης μετά από σφάλματα.
\item \textbf{Σειριακής σύνδεσης}: ο Ελεγχόμενος με Θυρίστορ Πυκνωτής Σειράς (Thyristor Controlled
      Series Capacitor, TCSC) και ο SSSC. Μεταβάλλουν τη φαινόμενη αντίδραση της γραμμής,
      ανακατευθύνοντας έτσι τη ροή ισχύος και αυξάνοντας τη μεταφορική ικανότητα.
\item \textbf{Συνδυασμένης λειτουργίας}: ο Ενοποιημένος Ελεγκτής Ροής Ισχύος (Unified Power Flow
      Controller, UPFC), ο οποίος ελέγχει ταυτόχρονα και ανεξάρτητα τάση, ενεργό και άεργο ροή.
\end{itemize}

Σε συνδυασμό με τους συνδέσμους συνεχούς ρεύματος υψηλής τάσης (HVDC), οι διατάξεις FACTS
αποτελούν τα βασικά εργαλεία του Διαχειριστή για τη ρύθμιση τάσης και την ανακούφιση συμφορήσεων
στο επίπεδο μεταφοράς. Τα αντίστοιχα εργαλεία στο επίπεδο διανομής εξετάζονται στα Κεφάλαια 5 και 6.

\subsection{Εκτίμηση κατάστασης}

\subsubsection*{Το πρόβλημα και το μοντέλο μέτρησης}

Οι μετρήσεις του SCADA είναι πλεονάζουσες, θορυβώδεις και ετερογενείς. Η \textbf{εκτίμηση
κατάστασης} (state estimation) είναι η διαδικασία που τις μετατρέπει σε μία συνεπή, πλήρη και
στατιστικά βέλτιστη εικόνα του δικτύου. Εισήχθη στα συστήματα ηλεκτρικής ενέργειας από τους
Schweppe και Wildes το 1970 [1, 2, 3] και αποτελεί έκτοτε τη θεμελιώδη λειτουργία κάθε EMS [5].

Ως \emph{κατάσταση} ορίζεται το διάνυσμα
$\mathbf{x} = [\boldsymbol{\theta}^{T}, \mathbf{V}^{T}]^{T}$ των γωνιών και των μέτρων των τάσεων
όλων των ζυγών. Για δίκτυο $N$ ζυγών, με έναν ζυγό αναφοράς, ισχύει $n = 2N-1$. Το μοντέλο μέτρησης
είναι

\begin{equation}
\mathbf{z} = \mathbf{h}(\mathbf{x}) + \mathbf{e}
\label{eq:measmodel}
\end{equation}

όπου $\mathbf{z} \in \mathbb{R}^{m}$ το διάνυσμα των μετρήσεων, $\mathbf{h}(\cdot)$ οι μη γραμμικές
συναρτήσεις που συνδέουν την κατάσταση με τα μετρούμενα μεγέθη (ροές και εγχύσεις ισχύος, μέτρα
τάσεων), και $\mathbf{e}$ ο θόρυβος μέτρησης, ο οποίος θεωρείται μηδενικής μέσης τιμής με
διαγώνιο πίνακα συνδιακύμανσης $\mathbf{R} = \mathrm{diag}(\sigma_1^2, \ldots, \sigma_m^2)$.

Καθοριστικό μέγεθος είναι ο \textbf{λόγος πλεονασμού} $\eta = m/n$, ο οποίος σε πραγματικά
συστήματα κυμαίνεται τυπικά μεταξύ 1,7 και 3. Ο πλεονασμός είναι εκείνος που επιτρέπει τόσο τη
μείωση του θορύβου όσο και την ανίχνευση εσφαλμένων μετρήσεων.

\subsubsection*{Σταθμισμένα ελάχιστα τετράγωνα}

Η καθιερωμένη διατύπωση είναι εκείνη των \textbf{σταθμισμένων ελαχίστων τετραγώνων} (Weighted Least
Squares, WLS), όπου ελαχιστοποιείται το σταθμισμένο άθροισμα των τετραγώνων των υπολοίπων:

\begin{equation}
J(\mathbf{x}) = \sum_{i=1}^{m} \frac{\left(z_i - h_i(\mathbf{x})\right)^2}{\sigma_i^2}
             = \left[\mathbf{z}-\mathbf{h}(\mathbf{x})\right]^{T} \mathbf{R}^{-1}
               \left[\mathbf{z}-\mathbf{h}(\mathbf{x})\right]
\label{eq:wls}
\end{equation}

Η στάθμιση με $1/\sigma_i^2$ σημαίνει ότι οι ακριβέστερες μετρήσεις βαρύνουν περισσότερο. Επειδή η
$\mathbf{h}$ είναι μη γραμμική, η λύση προκύπτει επαναληπτικά με τη μέθοδο Gauss--Newton. Με
$\mathbf{H}(\mathbf{x}) = \partial \mathbf{h}/\partial \mathbf{x}$ τον Ιακωβιανό πίνακα, σε κάθε
επανάληψη $k$ επιλύεται το σύστημα των κανονικών εξισώσεων

\begin{align}
\mathbf{G}(\mathbf{x}^{k}) \, \Delta \mathbf{x}^{k}
  &= \mathbf{H}^{T}(\mathbf{x}^{k}) \, \mathbf{R}^{-1}
     \left[\mathbf{z} - \mathbf{h}(\mathbf{x}^{k})\right] \\
\mathbf{G}(\mathbf{x}^{k})
  &= \mathbf{H}^{T}(\mathbf{x}^{k}) \, \mathbf{R}^{-1} \mathbf{H}(\mathbf{x}^{k})
\label{eq:gain}
\end{align}

και ενημερώνεται η εκτίμηση ως $\mathbf{x}^{k+1} = \mathbf{x}^{k} + \Delta \mathbf{x}^{k}$, έως
ότου $\max|\Delta \mathbf{x}^{k}|$ πέσει κάτω από ένα κατώφλι σύγκλισης. Ο πίνακας $\mathbf{G}$
ονομάζεται \emph{πίνακας κέρδους} (gain matrix) και είναι αραιός και συμμετρικός· η αποδοτική
παραγοντοποίησή του με τεχνικές αραιών πινάκων είναι αυτή που καθιστά το πρόβλημα επιλύσιμο σε
πραγματικό χρόνο για δίκτυα χιλιάδων ζυγών [4, 6].

\subsubsection*{Παρατηρησιμότητα}

Η επίλυση προϋποθέτει ότι ο $\mathbf{G}$ είναι αντιστρέψιμος, δηλαδή ότι το σύστημα είναι
\textbf{παρατηρήσιμο}: οι διαθέσιμες μετρήσεις αρκούν για τον μονοσήμαντο προσδιορισμό ολόκληρης
της κατάστασης. Ο έλεγχος γίνεται είτε τοπολογικά (αναζήτηση δένδρου ζεύξης καλυμμένου από
μετρήσεις) είτε αριθμητικά (έλεγχος βαθμού του $\mathbf{H}$). Όταν το σύστημα δεν είναι πλήρως
παρατηρήσιμο, εντοπίζονται οι \emph{παρατηρήσιμες νησίδες} και η κατάσταση εκτιμάται μόνο εντός
αυτών, με χρήση ψευδομετρήσεων (pseudo-measurements) για τη γεφύρωση των κενών [4, 6].

\subsubsection*{Ανίχνευση και αναγνώριση εσφαλμένων μετρήσεων}

Ο πλεονασμός επιτρέπει τον έλεγχο της συνέπειας των μετρήσεων. Η κλασική διαδικασία έχει δύο στάδια:

\begin{itemize}
\item \textbf{Ανίχνευση}: η αντικειμενική συνάρτηση $J(\hat{\mathbf{x}})$ στη λύση ακολουθεί
      κατανομή $\chi^2$ με $m-n$ βαθμούς ελευθερίας. Αν υπερβεί το κατώφλι
      $\chi^2_{m-n,\,\alpha}$ για δεδομένο επίπεδο εμπιστοσύνης, υπάρχει ένδειξη εσφαλμένων δεδομένων.
\item \textbf{Αναγνώριση}: υπολογίζονται τα κανονικοποιημένα υπόλοιπα
      $r_i^{N} = |r_i| / \sqrt{\Omega_{ii}}$, όπου $\mathbf{r} = \mathbf{z}-\mathbf{h}(\hat{\mathbf{x}})$
      και $\boldsymbol{\Omega} = \mathbf{R} - \mathbf{H}\mathbf{G}^{-1}\mathbf{H}^{T}$ ο πίνακας
      συνδιακύμανσης των υπολοίπων. Η μέτρηση με το μέγιστο $r_i^{N}$ (κριτήριο Largest Normalized
      Residual) απορρίπτεται και η εκτίμηση επαναλαμβάνεται [4, 6].
\end{itemize}

Δύο περιορισμοί της μεθόδου αξίζει να τονιστούν. Πρώτον, οι \emph{κρίσιμες μετρήσεις} — εκείνες
των οποίων η αφαίρεση καθιστά το σύστημα μη παρατηρήσιμο — έχουν εξ ορισμού μηδενικό υπόλοιπο και
επομένως ένα σφάλμα τους είναι μη ανιχνεύσιμο. Δεύτερον, όπως έδειξαν οι Liu, Ning και Reiter [17],
ένας επιτιθέμενος που γνωρίζει την τοπολογία και τις παραμέτρους του δικτύου μπορεί να κατασκευάσει
διάνυσμα επίθεσης $\mathbf{a} = \mathbf{H}\mathbf{c}$ το οποίο αφήνει το υπόλοιπο αμετάβλητο και
συνεπώς διαφεύγει πλήρως των παραπάνω ελέγχων. Οι επιθέσεις αυτές, γνωστές ως \emph{επιθέσεις
έγχυσης ψευδών δεδομένων} (False Data Injection Attacks, FDIA), εξετάζονται αναλυτικά στο Κεφάλαιο 8.

\subsubsection*{Ενσωμάτωση μετρήσεων PMU}

Η προσθήκη μετρήσεων PMU μεταβάλλει ουσιωδώς το πρόβλημα. Καθώς η PMU μετρά απευθείας τη γωνία
τάσης, η σχέση μεταξύ μετρήσεων και κατάστασης γίνεται \emph{γραμμική} όταν εκφραστεί σε ορθογώνιες
συντεταγμένες. Αν το σύστημα είναι πλήρως παρατηρήσιμο μόνο από PMU, η \eqref{eq:measmodel}
εκφυλίζεται σε $\mathbf{z} = \mathbf{H}\mathbf{x} + \mathbf{e}$ και η λύση δίνεται σε ένα βήμα,
χωρίς επαναλήψεις και χωρίς προβλήματα σύγκλισης:

\begin{equation}
\hat{\mathbf{x}} = \left(\mathbf{H}^{T}\mathbf{R}^{-1}\mathbf{H}\right)^{-1}
                    \mathbf{H}^{T}\mathbf{R}^{-1}\mathbf{z}
\label{eq:linse}
\end{equation}

Στην πράξη, η πλήρης κάλυψη με PMU σπανίζει για λόγους κόστους, οπότε χρησιμοποιούνται
\textbf{υβριδικοί εκτιμητές} που συνδυάζουν συμβατικές μετρήσεις SCADA με συγχρονοφασιθέτες [9].

\subsection{Στάδια λειτουργίας: από τον προγραμματισμό στον πραγματικό χρόνο}

Η λειτουργία του συστήματος δεν αποφασίζεται σε μία στιγμή, αλλά μέσα από διαδοχικές χρονικές
κλίμακες, καθεμιά από τις οποίες διορθώνει τα σφάλματα πρόβλεψης της προηγούμενης. Στο ευρωπαϊκό
μοντέλο αγοράς (Target Model) οι κλίμακες αυτές αντιστοιχούν σε διακριτές αγορές [13].

\begin{itemize}
\item \textbf{Μακροπρόθεσμος και προθεσμιακός ορίζοντας} (έτη έως εβδομάδες). Μελέτες επάρκειας
      ισχύος, συντονισμός προγραμμάτων συντήρησης γραμμών και μονάδων, εκτίμηση των αναγκών σε
      εφεδρείες και προθεσμιακά προϊόντα αντιστάθμισης κινδύνου.
\item \textbf{Αγορά Επόμενης Ημέρας} (day-ahead). Είναι το κύριο στάδιο προγραμματισμού. Οι
      συμμετέχοντες υποβάλλουν προσφορές για κάθε αγοραία περίοδο της επόμενης ημέρας και η αγορά
      εκκαθαρίζεται με ενιαία σύζευξη σε ευρωπαϊκό επίπεδο. Ο Διαχειριστής ελέγχει στη συνέχεια αν
      το προκύπτον πρόγραμμα είναι τεχνικά εφικτό.
\item \textbf{Ενδοημερήσια Αγορά} (intra-day). Επιτρέπει τη διόρθωση των θέσεων καθώς πλησιάζει η
      ώρα παράδοσης και βελτιώνονται οι προβλέψεις ζήτησης και — κυρίως — παραγωγής από ανανεώσιμες
      πηγές. Έχει ιδιαίτερη σημασία σε συστήματα με υψηλή διείσδυση αιολικής και φωτοβολταϊκής
      παραγωγής, όπου το σφάλμα πρόβλεψης μειώνεται σημαντικά όσο μικραίνει ο ορίζοντας.
\item \textbf{Πραγματικός χρόνος και εξισορρόπηση} (balancing). Ο Διαχειριστής ενεργοποιεί
      προσφορές ενέργειας εξισορρόπησης ώστε να καλύψει τις αποκλίσεις που απομένουν, να
      αντιμετωπίσει διαταραχές και να διαχειριστεί συμφορήσεις.
\end{itemize}

Στην Ελλάδα το μοντέλο αυτό τέθηκε σε ισχύ την 1η Νοεμβρίου 2020. Η Αγορά Επόμενης Ημέρας και η
Ενδοημερήσια Αγορά λειτουργούν από το Ελληνικό Χρηματιστήριο Ενέργειας (ΕΧΕ) [20], ενώ την
\textbf{Αγορά Εξισορρόπησης} λειτουργεί ο ΑΔΜΗΕ [16]. Ο τελευταίος εκτελεί τη Διαδικασία
Ολοκληρωμένου Προγραμματισμού και ενεργοποιεί σε πραγματικό χρόνο τόσο ενέργεια εξισορρόπησης όσο
και τις εφεδρείες που περιγράφονται στη συνέχεια.

\subsection{Απόκριση συχνότητας και εφεδρείες}

Η συχνότητα είναι ο πλέον άμεσος δείκτης του ισοζυγίου ισχύος: παραμένει σταθερή όσο παραγωγή και
ζήτηση εξισορροπούνται και αποκλίνει μόλις εμφανιστεί ανισοζύγιο. Αμέσως μετά από μια διαταραχή
ισχύος $\Delta P$, ο αρχικός ρυθμός μεταβολής της συχνότητας (Rate of Change of Frequency, ROCOF)
καθορίζεται αποκλειστικά από την \emph{αδράνεια} των συγχρονισμένων στρεφόμενων μαζών:

\begin{equation}
\frac{\mathrm{d}f}{\mathrm{d}t} = \frac{f_0 \, \Delta P}{2 H S_n}
\label{eq:rocof}
\end{equation}

όπου $f_0$ η ονομαστική συχνότητα, $H$ η σταθερά αδράνειας του συστήματος και $S_n$ η συνολική
ονομαστική φαινόμενη ισχύς των συγχρονισμένων μονάδων [8]. Η σχέση \eqref{eq:rocof} εξηγεί γιατί η
αντικατάσταση σύγχρονων γεννητριών από μονάδες συνδεδεμένες μέσω μετατροπέων — οι οποίες δεν
προσφέρουν εγγενή αδράνεια — οδηγεί σε ταχύτερες αποκλίσεις συχνότητας και καθιστά την απόκριση
απαιτητικότερη.

Η ανάκτηση της συχνότητας γίνεται μέσω τεσσάρων διαδοχικών επιπέδων εφεδρείας, τα οποία ορίζονται
στον SO GL [12] και εξειδικεύονται για την Ηπειρωτική Ευρώπη στη Συμφωνία Πλαισίου Σύγχρονης
Περιοχής (SAFA) [14]:

\begin{itemize}
\item \textbf{FCR} — Εφεδρεία Διατήρησης Συχνότητας (πρωτεύουσα ρύθμιση). Ενεργοποιείται αυτόματα
      και τοπικά από τους ρυθμιστές στροφών των μονάδων, ανάλογα προς την απόκλιση συχνότητας
      (στατισμός). \emph{Σταθεροποιεί} τη συχνότητα σε νέα μόνιμη τιμή, χωρίς να την επαναφέρει
      στα 50 Hz. Ενεργοποιείται από όλες τις μονάδες της σύγχρονης περιοχής από κοινού.
\item \textbf{aFRR} — Αυτόματη Εφεδρεία Αποκατάστασης Συχνότητας (δευτερεύουσα ρύθμιση).
      Ενεργοποιείται αυτόματα από τον κεντρικό ελεγκτή του κάθε Διαχειριστή, ο οποίος επιδιώκει τον
      μηδενισμό του Σφάλματος Ελέγχου Περιοχής (Area Control Error, ACE):
      $\mathrm{ACE} = \Delta P + K\,\Delta f$, όπου $\Delta P$ η απόκλιση των διασυνοριακών
      ανταλλαγών από το πρόγραμμα και $K$ ο συντελεστής της περιοχής. Επαναφέρει τη συχνότητα στα
      50 Hz και απελευθερώνει την FCR.
\item \textbf{mFRR} — Χειροκίνητη Εφεδρεία Αποκατάστασης Συχνότητας (τριτεύουσα ρύθμιση).
      Ενεργοποιείται με εντολή του Διαχειριστή και αντικαθιστά την aFRR σε παρατεταμένα ανισοζύγια.
\item \textbf{RR} — Εφεδρεία Αντικατάστασης. Αποκαθιστά τα επίπεδα των προηγούμενων εφεδρειών ώστε
      το σύστημα να είναι έτοιμο για επόμενη διαταραχή.
\end{itemize}

\begin{table}[!htbp]
\centering
\begin{tabular}{lllc}
\hline
Εφεδρεία & Ονομασία & Ενεργοποίηση & Πλήρης ενεργοποίηση \\
\hline
FCR  & Πρωτεύουσα   & αυτόματη, τοπική  & 30 s \\
aFRR & Δευτερεύουσα & αυτόματη, κεντρική & 5 min \\
mFRR & Τριτεύουσα   & χειροκίνητη        & 12,5 min \\
RR   & Αντικατάστασης & χειροκίνητη      & $>$ 15 min \\
\hline
\end{tabular}
\caption{Επίπεδα εφεδρείας συχνότητας κατά τον SO GL. Οι χρόνοι πλήρους ενεργοποίησης αφορούν
τη σύγχρονη περιοχή της Ηπειρωτικής Ευρώπης [12, 14].}
\label{tab:reserves}
\end{table}

Οι ποσοτικές απαιτήσεις για την Ηπειρωτική Ευρώπη είναι συγκεκριμένες [12, 14]:

\begin{itemize}
\item Το \textbf{συμβάν αναφοράς} για τη διαστασιολόγηση της FCR είναι απώλεια 3\,000 MW, και
      αντίστοιχα 3\,000 MW στην αντίθετη κατεύθυνση. Η ελάχιστη διαθέσιμη FCR στη σύγχρονη περιοχή
      είναι ανά πάσα στιγμή 3\,000 MW.
\item Η μέγιστη απόκλιση συχνότητας μόνιμης κατάστασης για το συμβάν αναφοράς είναι
      \textbf{200 mHz}, οπότε ο συντελεστής αναλογίας μεταξύ FCR και απόκλισης συχνότητας
      προκύπτει ίσος με 15\,000 MW/Hz.
\item Το \textbf{τυπικό εύρος συχνότητας} είναι $\pm 50$ mHz και δεν επιτρέπεται να παραβιάζεται
      για περισσότερα από 15\,000 λεπτά ανά έτος. Η μέγιστη στιγμιαία απόκλιση ορίζεται σε
      800 mHz.
\item Απόκλιση συχνότητας μεγαλύτερη των 200 mHz οδηγεί σε κατάσταση \emph{έκτακτης ανάγκης}.
\end{itemize}

\subsection{Σύνοψη και σύνδεση με τα επόμενα κεφάλαια}

Η λειτουργία του ΣΜΗΕ στηρίζεται σε έναν κλειστό βρόχο: η υποδομή εποπτείας (SCADA, PMU) παράγει
μετρήσεις, η εκτίμηση κατάστασης τις μετατρέπει σε αξιόπιστη εικόνα του δικτύου, τα εργαλεία
ανάλυσης ασφάλειας αξιολογούν την εικόνα αυτή, και οι αποφάσεις υλοποιούνται μέσω των αγορών και
των διατάξεων ελέγχου σε κλίμακες που εκτείνονται από τα δευτερόλεπτα της πρωτεύουσας ρύθμισης έως
τους μήνες του προγραμματισμού.

Οι υπηρεσίες που στηρίζουν αυτόν τον βρόχο ονομάζονται \emph{επικουρικές υπηρεσίες} και
διακρίνονται σε σχετικές και μη σχετικές με τη συχνότητα· τις εξετάζουν αντίστοιχα τα Κεφάλαια 3
και 2. Το αντίστοιχο πρόβλημα στο επίπεδο της διανομής, όπου η παρατηρησιμότητα είναι πολύ
χαμηλότερη, εξετάζεται στο Κεφάλαιο 4, ενώ η ασφάλεια της ίδιας της υποδομής μέτρησης και ελέγχου
αποτελεί το αντικείμενο του Κεφαλαίου 8.

\section{Παραδείγματα}
\label{kef-01-paradeigmata}
% TODO: παραδείγματα και λυμένες ασκήσεις

\section{Βιβλιογραφία}
\label{kef-01-vivliografia}

\begin{enumerate}

\item F. C. Schweppe and J. Wildes, ``Power System Static-State Estimation, Part I: Exact Model,''
      \textit{IEEE Transactions on Power Apparatus and Systems}, vol. PAS-89, no. 1, pp. 120--125,
      Jan. 1970. doi: \texttt{10.1109/TPAS.1970.292678}

\item F. C. Schweppe and D. B. Rom, ``Power System Static-State Estimation, Part II: Approximate
      Model,'' \textit{IEEE Transactions on Power Apparatus and Systems}, vol. PAS-89, no. 1,
      pp. 125--130, Jan. 1970.

\item F. C. Schweppe, ``Power System Static-State Estimation, Part III: Implementation,''
      \textit{IEEE Transactions on Power Apparatus and Systems}, vol. PAS-89, no. 1, pp. 130--135,
      Jan. 1970.

\item A. Monticelli, \textit{State Estimation in Electric Power Systems: A Generalized Approach}.
      Boston, MA: Kluwer Academic Publishers, 1999.

\item A. Monticelli, ``Electric Power System State Estimation,'' \textit{Proceedings of the IEEE},
      vol. 88, no. 2, pp. 262--282, Feb. 2000.
      doi: \texttt{10.1109/5.824004}

\item A. Abur and A. Gómez Expósito, \textit{Power System State Estimation: Theory and
      Implementation}. New York: Marcel Dekker, 2004.

\item A. J. Wood, B. F. Wollenberg and G. B. Sheblé, \textit{Power Generation, Operation, and
      Control}, 3rd ed. Hoboken, NJ: Wiley, 2014.

\item P. Kundur, \textit{Power System Stability and Control}. New York: McGraw-Hill, 1994.

\item A. G. Phadke and J. S. Thorp, \textit{Synchronized Phasor Measurements and Their
      Applications}, 2nd ed. Cham: Springer, 2017.

\item N. G. Hingorani and L. Gyugyi, \textit{Understanding FACTS: Concepts and Technology of
      Flexible AC Transmission Systems}. New York: IEEE Press, 2000.

\item IEC/IEEE 60255-118-1:2018, \textit{Measuring Relays and Protection Equipment --- Part 118-1:
      Synchrophasor for Power Systems --- Measurements}. IEC/IEEE, 2018.

\item Commission Regulation (EU) 2017/1485 of 2 August 2017 establishing a guideline on electricity
      transmission system operation (System Operation Guideline, SO GL). \textit{Official Journal
      of the European Union}, L 220, 25.8.2017.
      \href{https://eur-lex.europa.eu/eli/reg/2017/1485/oj}{eur-lex.europa.eu/eli/reg/2017/1485/oj}

\item Commission Regulation (EU) 2017/2195 of 23 November 2017 establishing a guideline on
      electricity balancing (Electricity Balancing Guideline, EB GL). \textit{Official Journal of
      the European Union}, L 312, 28.11.2017.
      \href{https://eur-lex.europa.eu/eli/reg/2017/2195/oj}{eur-lex.europa.eu/eli/reg/2017/2195/oj}

\item ENTSO-E, \textit{Synchronous Area Framework Agreement for Regional Group Continental Europe
      --- Policy 1: Load-Frequency Control and Reserves}, 2019.
      \href{https://www.entsoe.eu/}{entsoe.eu}

\item ΑΔΜΗΕ Α.Ε., \textit{Κώδικας Διαχείρισης του Ελληνικού Συστήματος Μεταφοράς Ηλεκτρικής
      Ενέργειας}. \href{https://www.admie.gr/}{admie.gr}

\item ΑΔΜΗΕ Α.Ε., \textit{Κανονισμός Αγοράς Εξισορρόπησης}. \href{https://www.admie.gr/}{admie.gr}

\item Y. Liu, P. Ning and M. K. Reiter, ``False Data Injection Attacks against State Estimation in
      Electric Power Grids,'' \textit{ACM Transactions on Information and System Security},
      vol. 14, no. 1, art. 13, May 2011.
      doi: \texttt{10.1145/1952982.1952995}

\item IEC 61850, \textit{Communication Networks and Systems for Power Utility Automation}, και
      IEC 60870-5-104, \textit{Telecontrol Equipment and Systems --- Part 5-104: Network Access for
      IEC 60870-5-101 Using Standard Transport Profiles}. IEC.

\item T. E. Dy Liacco, ``The Adaptive Reliability Control System,'' \textit{IEEE Transactions on
      Power Apparatus and Systems}, vol. PAS-86, no. 5, pp. 517--531, May 1967.
      doi: \texttt{10.1109/TPAS.1967.291728}

\item Ελληνικό Χρηματιστήριο Ενέργειας (ΕΧΕ), \textit{Κανονισμός Αγοράς Επόμενης Ημέρας και
      Ενδοημερήσιας Αγοράς}. \href{https://www.enexgroup.gr/}{enexgroup.gr}

\end{enumerate}
